\pagestyle{empty}
\tight
\begin{tblr}{width=\textwidth,colspec={|Q[c,m,1.9cm]|X[l,m]|},hlines,vlines,hspan=minimal,row{1}={bg=headblue},rowsep=1pt,colsep=3pt}
\SetCell{c,m}\hfil\logo\hfil & \textbf{POLITEKNIK NEGERI MALANG}\par\textbf{JURUSAN TEKNOLOGI INFORMASI}\par\textbf{PROGRAM STUDI: D4 TEKNIK INFORMATIKA} \\
\SetCell[c=2]{c} \textbf{RENCANA PEMBELAJARAN SEMESTER (RPS)} & \\
\end{tblr}
\begin{longtblr}{width=\textwidth,colspec={|Q[l,m,1.9cm]|X[0.85,l,m]|X[1.45,l,m]|X[0.75,l,m]|X[1.15,l,m]|},hlines,vlines,hspan=minimal,rowsep=1pt,colsep=2pt,row{1,3}={bg=headgray,font=\bfseries}}
MATA KULIAH & KODE & BOBOT (SKS)/jam & SEMESTER & TGL. PENYUSUNAN \\
Bahasa Inggris 2 & RTI253009 & 2 SKS/4 jam & 3 & 1 Juli 2025 \\
OTORISASI & Dosen Pengembang RPS & Koordinator RMK & \SetCell[c=2]{l} Ka PRODI &  \\
 & Farid Angga Pribadi, S.Kom., M.Kom. &  & \SetCell[c=2]{l}{\vspace{8mm}\par Dr. Ely Setyo Astuti, S.T., M.T.} &  \\
\SetCell[r=9]{l} Capaian Pembelajaran (CP) & \SetCell[c=4]{l,bg=headgray,font=\bfseries} Capaian Pembelajaran Lulusan Program Studi (CPL-Prodi) &  &  &  \\
 & CPL01 & \SetCell[c=3]{l} Mampu menerapkan nilai ketakwaan, kesadaran hukum, kedisiplinan, profesionalisme, etika profesi, dan pembelajaran sepanjang hayat, serta tanggap terhadap isu sosial dan perkembangan teknologi. &  &  \\
 & CPL03 & \SetCell[c=3]{l} Mampu mengelola tim dan diri sendiri secara profesional serta mengomunikasikan dan mempresentasikan gagasan maupun hasil kerja secara efektif baik lisan maupun tulisan. &  &  \\
 & \SetCell[c=4]{l,bg=headgray,font=\bfseries} Capaian Pembelajaran mata kuliah (CPL-MK) &  &  &  \\
 & CPMK0102 & \SetCell[c=3]{l} Mahasiswa mampu menunjukkan sikap profesional, pembelajaran sepanjang hayat, dan tanggung jawab dalam konteks akademik internasional. &  &  \\
 & CPMK0303 & \SetCell[c=3]{l} Mahasiswa mampu mengembangkan kemampuan komunikasi profesional dalam bahasa Inggris untuk keperluan akademik dan teknis. &  &  \\
 & \SetCell[c=4]{l,bg=headgray,font=\bfseries} Kemampuan akhir tiap tahapan belajar (Sub-CPMK) &  &  &  \\
 & SCPMK0102-02701 & \SetCell[c=3]{l} Mahasiswa mampu membaca dan memahami teks teknis bahasa Inggris dalam bidang TI serta mengekspresikan ide akademik secara tertulis. &  &  \\
 & SCPMK0303-02702 & \SetCell[c=3]{l} Mahasiswa mampu berkomunikasi lisan dalam bahasa Inggris pada konteks diskusi teknis, presentasi, dan wawancara simulasi. &  &  \\
\SetCell[r=2]{l} Penilaian & \SetCell[c=4]{l} *Nilai CPMK didapat dari agregasi nilai SUB-CPMK yang dibebankan &  &  &  \\
 & \SetCell[c=4]{l}{\tiny\begin{tblr}{width=\linewidth,colspec={|Q[c,m,0.1\linewidth]|Q[c,m,0.16\linewidth]|X[l,m]|X[l,m]|X[l,m]|X[l,m]|X[l,m]|X[l,m]|Q[c,m,0.08\linewidth]|},hlines,vlines,hspan=minimal,rowsep=1pt,colsep=0.7pt,row{1,2}={bg=headgray,font=\bfseries}}
Kode CPMK & Kode SCPMK & \SetCell[c=6]{c} Bobot per Bentuk Penilaian & & & & & & Total Bobot \\
 & & Tugas & Kuis 1 & UTS & Kuis 2 & PBL & UAS & \\
CPMK0102 & SCPMK0102-02701 & 15\%: tugas writing dan reading & 5\%: kosakata dan bacaan & 10\%: UTS comprehension & & & 10\%: UAS writing & 40\% \\
CPMK0303 & SCPMK0303-02702 & 15\%: tugas speaking dan presentasi & & 5\%: UTS komunikasi lisan & 5\%: kuis listening dan speaking & 22\%: presentasi teknis & 13\%: UAS presentasi akhir & 60\% \\
\SetCell[c=8]{r}\textbf{Total Per Penilaian} & & & & & & & & \textbf{100\%} \\
\end{tblr}} &  &  &  \\
Diskripsi Singkat Mata Kuliah & \SetCell[c=4]{l} Bahasa Inggris 2 mengembangkan kemampuan membaca teks teknis, menulis laporan akademik, dan berkomunikasi lisan dalam bahasa Inggris pada konteks profesional dan akademik bidang TI. &  &  &  \\
Bahan Kajian / Materi Pembelajaran & \SetCell[c=4]{l}\begin{enumerate}\item Reading: kosakata teknis IT, pemahaman teks teknis, dan membaca literatur\item Writing: paragraf akademik, laporan teknis, dan CV serta email profesional\item Listening: ceramah dan presentasi teknis\item Speaking: diskusi, presentasi teknis, dan wawancara simulasi\end{enumerate} &  &  &  \\
\SetCell[r=2]{l} Pustaka & \SetCell[c=4]{l} \textbf{Utama:}\begin{enumerate}\item Murphy, R. 2019. English Grammar in Use. Cambridge University Press.\item Glendinning, E. 2008. Professional English in Use IT. Cambridge University Press.\item Swan, M. 2016. Practical English Usage. Oxford University Press.\end{enumerate} &  &  &  \\
 & \SetCell[c=4]{l} \textbf{Pendukung:} Materi LMS, artikel teknis IT berbahasa Inggris, dan sumber daring terpercaya. &  &  &  \\
Media Pembelajaran & \SetCell[c=2]{l} \textbf{Software:} Browser, LMS, aplikasi audio/video, office suite. &  & \SetCell[c=2]{l} \textbf{Hardware:} Komputer/laptop, speaker, mikrofon, proyektor. &  \\
Nama Dosen Pengampu & \SetCell[c=4]{l} Tim Pengajar Program Studi D4 TI &  &  &  \\
Matakuliah Syarat & \SetCell[c=4]{l} - &  &  &  \\
\end{longtblr}
\begin{longtblr}{width=\textwidth,colspec={|Q[c,m,0.06\textwidth]|X[1.35,l,m]|X[1.08,l,m]|X[1.16,l,m]|X[0.7,l,m]|X[1.24,l,m]|X[1.06,l,m]|X[0.98,l,m]|Q[c,m,0.05\textwidth]|},hlines,vlines,hspan=minimal,rowhead=2,row{1,2}={bg=headgreen,font=\bfseries},rowsep=1pt,colsep=1.5pt}
Mgg.\par Ke & Kemampuan Akhir Yang Direncanakan (Sub-CP-MK) & Bahan kajian (Materi Pembelajaran) & Bentuk dan Metode Pembelajaran & Estimasi Waktu & Pengalaman Belajar Mahasiswa & Kriteria \& Bentuk Penilaian & Indikator Penilaian & Bobot Penilaian (\%) \\
(1) & (2) & (3) & (4) & (5) & (6) & (7) & (8) & (9) \\
1 & SCPMK0102-02701 & Reading technical texts: IT vocabulary and comprehension. & Modalitas: Blended Learning. Bentuk: Luring/Daring. Metode: Case Method, diskusi kelompok, presentasi, dan peer feedback. & 1 x 2 x 50' tatap muka; 1 x 2 x 50' tugas/praktik mandiri. & Mahasiswa membaca teks teknis IT, mengidentifikasi kosakata, dan menjawab pertanyaan pemahaman. & Bentuk penilaian: Tugas Reading. Mengacu rubrik RTM. & Ketepatan konsep; kualitas komunikasi; validasi dan dokumentasi. & 2.5 \\
2 & SCPMK0102-02701 & Academic writing: paragraph structure and essay basics. & Modalitas: Blended Learning. Bentuk: Luring/Daring. Metode: Case Method, diskusi kelompok, presentasi, dan peer feedback. & 1 x 2 x 50' tatap muka; 1 x 2 x 50' tugas/praktik mandiri. & Mahasiswa menulis paragraf akademik dengan struktur yang benar dan berlatih penulisan esai pendek. & Bentuk penilaian: Tugas Writing. Mengacu rubrik RTM. & Ketepatan konsep; kualitas komunikasi; validasi dan dokumentasi. & 2.5 \\
3 & SCPMK0102-02701 & Technical writing: report writing in English. & Modalitas: Blended Learning. Bentuk: Luring/Daring. Metode: Case Method, diskusi kelompok, presentasi, dan peer feedback. & 1 x 2 x 50' tatap muka; 1 x 2 x 50' tugas/praktik mandiri. & Mahasiswa menulis laporan teknis dalam bahasa Inggris sesuai format yang ditetapkan. & Bentuk penilaian: Tugas Writing. Mengacu rubrik RTM. & Ketepatan konsep; kualitas komunikasi; validasi dan dokumentasi. & 2.5 \\
4 & SCPMK0102-02701 & Kuis 1: reading and vocabulary. & Modalitas: Blended Learning. Bentuk: Luring/Daring. Metode: Case Method, diskusi kelompok, presentasi, dan peer feedback. & 1 x 2 x 50' tatap muka; 1 x 2 x 50' tugas/praktik mandiri. & Mahasiswa mengerjakan kuis untuk mengukur kemampuan membaca teks teknis dan penguasaan kosakata IT. & Bentuk penilaian: Kuis 1. Mengacu rubrik RTM. & Ketepatan konsep; kualitas komunikasi; validasi dan dokumentasi. & 5 \\
5 & SCPMK0303-02702 & Listening skills: lectures and technical presentations. & Modalitas: Blended Learning. Bentuk: Luring/Daring. Metode: Case Method, diskusi kelompok, presentasi, dan peer feedback. & 1 x 2 x 50' tatap muka; 1 x 2 x 50' tugas/praktik mandiri. & Mahasiswa mendengarkan presentasi teknis dan menjawab pertanyaan pemahaman mendengarkan. & Bentuk penilaian: Tugas Listening. Mengacu rubrik RTM. & Ketepatan konsep; kualitas komunikasi; validasi dan dokumentasi. & 2.5 \\
6 & SCPMK0303-02702 & Speaking: discussion and expressing opinions. & Modalitas: Blended Learning. Bentuk: Luring/Daring. Metode: Case Method, diskusi kelompok, presentasi, dan peer feedback. & 1 x 2 x 50' tatap muka; 1 x 2 x 50' tugas/praktik mandiri. & Mahasiswa berpartisipasi dalam diskusi kelompok teknis dan menyampaikan pendapat dalam bahasa Inggris. & Bentuk penilaian: Tugas Speaking. Mengacu rubrik RTM. & Ketepatan konsep; kualitas komunikasi; validasi dan dokumentasi. & 2.5 \\
7 & SCPMK0303-02702 & Technical presentation skills. & Modalitas: Blended Learning. Bentuk: Luring/Daring. Metode: Case Method, diskusi kelompok, presentasi, dan peer feedback. & 1 x 2 x 50' tatap muka; 1 x 2 x 50' tugas/praktik mandiri. & Mahasiswa berlatih menyampaikan presentasi teknis singkat dalam bahasa Inggris. & Bentuk penilaian: Tugas Presentasi. Mengacu rubrik RTM. & Ketepatan konsep; kualitas komunikasi; validasi dan dokumentasi. & 2.5 \\
8 & SCPMK0102-02701, SCPMK0303-02702 & UTS: reading comprehension and writing. & Modalitas: Blended Learning. Bentuk: Luring/Daring. Metode: Case Method, diskusi kelompok, presentasi, dan peer feedback. & 1 x 2 x 50' tatap muka; 1 x 2 x 50' tugas/praktik mandiri. & Mahasiswa mengerjakan ujian tengah semester yang mengukur kemampuan membaca, menulis, dan komunikasi lisan. & Bentuk penilaian: UTS. Mengacu rubrik RTM. & Ketepatan konsep; kualitas komunikasi; validasi dan dokumentasi. & 15 \\
9 & SCPMK0102-02701 & Research paper reading and summarizing. & Modalitas: Blended Learning. Bentuk: Luring/Daring. Metode: Case Method, diskusi kelompok, presentasi, dan peer feedback. & 1 x 2 x 50' tatap muka; 1 x 2 x 50' tugas/praktik mandiri. & Mahasiswa membaca artikel penelitian IT dan membuat ringkasan terstruktur dalam bahasa Inggris. & Bentuk penilaian: Tugas Reading. Mengacu rubrik RTM. & Ketepatan konsep; kualitas komunikasi; validasi dan dokumentasi. & 2.5 \\
10 & SCPMK0102-02701 & Critical reading: evaluating IT sources. & Modalitas: Blended Learning. Bentuk: Luring/Daring. Metode: Case Method, diskusi kelompok, presentasi, dan peer feedback. & 1 x 2 x 50' tatap muka; 1 x 2 x 50' tugas/praktik mandiri. & Mahasiswa mengevaluasi kredibilitas dan relevansi sumber teknis IT berbahasa Inggris. & Bentuk penilaian: Tugas Reading. Mengacu rubrik RTM. & Ketepatan konsep; kualitas komunikasi; validasi dan dokumentasi. & 2.5 \\
11 & SCPMK0303-02702 & Presentation delivery: slide design and delivery. & Modalitas: Blended Learning. Bentuk: Luring/Daring. Metode: Case Method, diskusi kelompok, presentasi, dan peer feedback. & 1 x 2 x 50' tatap muka; 1 x 2 x 50' tugas/praktik mandiri. & Mahasiswa merancang slide presentasi teknis dan berlatih delivery yang efektif dalam bahasa Inggris. & Bentuk penilaian: Tugas Presentasi. Mengacu rubrik RTM. & Ketepatan konsep; kualitas komunikasi; validasi dan dokumentasi. & 2.5 \\
12 & SCPMK0303-02702 & Kuis 2: listening and speaking skills. & Modalitas: Blended Learning. Bentuk: Luring/Daring. Metode: Case Method, diskusi kelompok, presentasi, dan peer feedback. & 1 x 2 x 50' tatap muka; 1 x 2 x 50' tugas/praktik mandiri. & Mahasiswa mengerjakan kuis untuk mengukur kemampuan mendengarkan dan berbicara dalam bahasa Inggris teknis. & Bentuk penilaian: Kuis 2. Mengacu rubrik RTM. & Ketepatan konsep; kualitas komunikasi; validasi dan dokumentasi. & 5 \\
13 & SCPMK0303-02702 & Group discussion: technical problem solving in English. & Modalitas: Blended Learning. Bentuk: Luring/Daring. Metode: Case Method, diskusi kelompok, presentasi, dan peer feedback. & 1 x 2 x 50' tatap muka; 1 x 2 x 50' tugas/praktik mandiri. & Mahasiswa berpartisipasi dalam diskusi kelompok untuk memecahkan masalah teknis IT dalam bahasa Inggris. & Bentuk penilaian: Tugas Speaking. Mengacu rubrik RTM. & Ketepatan konsep; kualitas komunikasi; validasi dan dokumentasi. & 2.5 \\
14 & SCPMK0102-02701 & CV writing and professional email in English. & Modalitas: Blended Learning. Bentuk: Luring/Daring. Metode: Case Method, diskusi kelompok, presentasi, dan peer feedback. & 1 x 2 x 50' tatap muka; 1 x 2 x 50' tugas/praktik mandiri. & Mahasiswa menulis CV dan email profesional dalam bahasa Inggris sesuai standar industri TI. & Bentuk penilaian: Tugas Writing. Mengacu rubrik RTM. & Ketepatan konsep; kualitas komunikasi; validasi dan dokumentasi. & 2.5 \\
15 & SCPMK0303-02702 & Mock interview in English. & Modalitas: Blended Learning. Bentuk: Luring/Daring. Metode: Case Method, diskusi kelompok, presentasi, dan peer feedback. & 1 x 2 x 50' tatap muka; 1 x 2 x 50' tugas/praktik mandiri. & Mahasiswa berlatih wawancara kerja simulasi dalam bahasa Inggris dengan skenario posisi IT. & Bentuk penilaian: Tugas Speaking. Mengacu rubrik RTM. & Ketepatan konsep; kualitas komunikasi; validasi dan dokumentasi. & 2.5 \\
16 & SCPMK0303-02702 & PBL: English technical presentation project. & Modalitas: Blended Learning. Bentuk: Luring/Daring. Metode: Case Method, diskusi kelompok, presentasi, dan peer feedback. & 1 x 2 x 50' tatap muka; 1 x 2 x 50' tugas/praktik mandiri. & Mahasiswa menyajikan proyek presentasi teknis dalam bahasa Inggris tentang topik TI pilihan secara menyeluruh. & Bentuk penilaian: PBL Presentasi Teknis. Mengacu rubrik RTM. & Ketepatan konsep; kualitas komunikasi; validasi dan dokumentasi. & 22 \\
17 & SCPMK0102-02701, SCPMK0303-02702 & UAS: final presentation and writing test. & Modalitas: Blended Learning. Bentuk: Luring/Daring. Metode: Case Method, diskusi kelompok, presentasi, dan peer feedback. & 1 x 2 x 50' tatap muka; 1 x 2 x 50' tugas/praktik mandiri. & Mahasiswa mengerjakan ujian akhir semester yang mencakup tes menulis dan presentasi final dalam bahasa Inggris. & Bentuk penilaian: UAS. Mengacu rubrik RTM. & Ketepatan konsep; kualitas komunikasi; validasi dan dokumentasi. & 23 \\
\end{longtblr}
